\todo{Chương 1 chỉ có khung: nhóm tự viết văn bản từ các ý chính bên dưới, trích dẫn nguồn đã đọc. Các trích dẫn dưới đây là gợi ý trong refs.bib; kiểm tra lại trang/phiên bản trước khi dùng. Tham chiếu nội dung môn học: giáo trình và slide chương 1, 2, 5, 6.}

\section{Giới thiệu về kho dữ liệu}
\begin{itemize}
  \item Khái niệm kho dữ liệu; định nghĩa của Inmon: hướng chủ đề, tích hợp, biến thiên theo thời gian, không biến động \cite{inmon2005}; cách nhìn của Kimball \cite{kimball2013}.
  \item Vì sao doanh nghiệp cần kho dữ liệu: hệ thống nguồn rời rạc, dữ liệu không nhất quán, nhu cầu phân tích lịch sử.
  \item Đặc tính của kho dữ liệu và lợi ích (một nguồn sự thật, định nghĩa KPI thống nhất).
\end{itemize}

\section{OLTP và OLAP}
\begin{itemize}
  \item So sánh hệ thống giao dịch (OLTP) và hệ thống phân tích (OLAP): mục đích, mô hình dữ liệu (chuẩn hóa so với chiều), kiểu truy vấn, khối lượng, người dùng \cite{chaudhuri1997}.
  \item Ví dụ minh họa lấy từ dự án: 9 bảng chuẩn hóa của hệ thống bán hàng so với 1 fact và các chiều (liên hệ Chương 3).
\end{itemize}

\section{Kiến trúc kho dữ liệu}
\begin{itemize}
  \item Kiến trúc Kimball: bus matrix, conformed dimension, data mart theo quy trình nghiệp vụ \cite{kimball2013}.
  \item Kiến trúc Inmon: kho dữ liệu doanh nghiệp (3NF) và data mart phụ thuộc \cite{inmon2005}.
  \item Kiến trúc lai, vai trò của staging và ODS; so sánh ưu nhược điểm.
  \item Các lớp dữ liệu: nguồn, staging, ODS, kho, data mart, BI (liên hệ kiến trúc nhóm chọn ở mục 3.3.2).
\end{itemize}

\section{Mô hình dữ liệu đa chiều}
\begin{itemize}
  \item Fact, dimension, measure, grain; lược đồ sao, bông tuyết, chòm sao (constellation).
  \item Ba loại bảng fact: transaction, periodic snapshot, accumulating snapshot \cite{kimball2013}.
  \item Khóa thay thế, dòng Unknown, degenerate dimension, conformed dimension, junk dimension, role-playing dimension.
  \item Slowly changing dimension: loại 1, 2, 3 và khi nào dùng.
  \item Phân cấp (hierarchy) và cách dùng trong roll-up / drill-down.
\end{itemize}

\section{Quy trình ETL}
\begin{itemize}
  \item Extract, Transform, Load; nạp đầy đủ và nạp tăng dần (incremental), watermark, upsert.
  \item Làm sạch dữ liệu, kiểm soát chất lượng, ghi nhận dòng lỗi.
  \item Công cụ ETL phổ biến và lựa chọn của nhóm (Python điều phối, biến đổi bằng SQL procedure).
\end{itemize}

\section{OLAP: MOLAP, ROLAP, HOLAP}
\begin{itemize}
  \item Các thao tác OLAP: roll-up, drill-down, slice, dice, pivot \cite{chaudhuri1997,codd1993}.
  \item \textbf{So sánh MOLAP / ROLAP / HOLAP} (lưu trữ, tốc độ truy vấn, khả năng mở rộng, chi phí xử lý khối, độ tươi của dữ liệu, công cụ điển hình): nhóm lập bảng so sánh và trích nguồn.
  \item \textbf{Vì sao nhóm chọn ROLAP}: dùng PostgreSQL (không có khối đa chiều như SSAS), dữ liệu cỡ vài trăm nghìn dòng fact nên SQL trên lược đồ sao trả lời trong vòng một giây, kho luôn cập nhật sau mỗi lần nạp tăng dần mà không phải xử lý lại khối, Power BI đóng vai trò lớp ngữ nghĩa; đổi lại cần materialized view cho truy vấn nặng (liên hệ mục 3.4.4 và 3.5.1).
\end{itemize}

\section{Data mart và kho dữ liệu trong kinh doanh}
\begin{itemize}
  \item Khái niệm data mart, độc lập và phụ thuộc; tổng hợp sẵn bằng view / materialized view.
  \item Ứng dụng kho dữ liệu trong bán lẻ, thương mại điện tử, quản lý đơn hàng (giáo trình chương 7).
\end{itemize}
